\section{E2E convites, autenticação e notificações}
\label{sec:e2e-convites-auth}

Esta seção registra a primeira suíte E2E permanente com Playwright Test,
derivada da Seção~9 (fluxos), da Seção~8 (UI e estados) e da Seção~7
(resultados e regras). O objetivo não é cobertura numérica, mas provar a
vertical que recentemente revelou bugs que testes isolados não detectaram:
autenticação, emissão de convite, callable, Auth/Firestore, notificação interna
ou OOB, acesso a \texttt{/convite}, aceite/rejeição, matrícula, atualização de
custom claims e autorização real após o aceite.

\subsection{Infraestrutura}

\begin{itemize}
  \item \texttt{frontend/playwright.config.ts} --- \texttt{testDir} em
    \texttt{./e2e/specs}, \texttt{workers: 1}, \texttt{webServer} para o Next.js
    em modo \texttt{dev} (obrigatório para a conexão aos emuladores) e
    \texttt{globalSetup} que reinicia e semeia o estado canônico.
  \item \texttt{frontend/e2e/} --- \texttt{fixtures.ts} (guarda de console/page),
    \texttt{global-setup.ts}, \texttt{orchestrator.mjs}, \texttt{helpers/} e
    \texttt{specs/}.
  \item \texttt{frontend/e2e/helpers/browsers.ts} --- resolve a store path
    \texttt{*-playwright-browsers} quando \texttt{PLAYWRIGHT\_BROWSERS\_PATH} não
    está definido; nenhum browser é instalado.
  \item \texttt{frontend/e2e/global-setup.ts} --- aborta com mensagem clara se
    Auth (9099), Firestore (8080) ou Functions (5001) não estiverem disponíveis.
    Sem fallback para Firebase real.
\end{itemize}

Comandos:

\begin{verbatim}
cd frontend
npm run test:e2e:list        # lista cenários
npm run test:e2e             # exige emuladores já ativos; reseta+semeia via globalSetup
npm run test:e2e:emulators   # sobe emuladores, roda a suíte e derruba tudo
\end{verbatim}

O seed usado é exclusivamente o canônico \texttt{functions/scripts/seed.ts}
(\texttt{--reset}). Nenhum segundo universo de fixtures foi criado.

\subsection{Cenários}

\begin{longtable}{@{}p{0.13\textwidth}p{0.44\textwidth}p{0.35\textwidth}@{}}
\toprule
ID & Cenário & Camadas \\
\midrule
\endhead
E2E-001 & Login por e-mail/senha, claims de papel no header e roteamento de
  rota protegida & Auth Emulator, SDK, AuthContext, ProtectedRoute \\
E2E-002 & Professor convida aluno existente pela UI sem campos opcionais &
  UI, callable \texttt{convidarAluno}, Firestore \\
E2E-003 & Inbox do destinatário mostra o convite; terceiro não vê nada e não há
  \texttt{permission-denied} & Firestore Rules, NotificacoesDropdown \\
E2E-004/E2E-008 & Aluno existente aceita via notificação, sem campo de
  matrícula & UI, callable \texttt{aceitarConviteAluno}, vínculo, notificação \\
E2E-005 & Convite externo provisiona Auth e preserva \texttt{continueUrl}/\texttt{token}
  do OOB real & Identity Toolkit \texttt{sendOobCode}, OOB, \texttt{/convite} \\
E2E-007 & Matrícula obrigatória quando ausente, com preservação de zeros &
  projeção \texttt{matricula\_necessaria}, callable, Firestore \\
E2E-006 & E-mail não verificado não aceita o convite & barreira \texttt{emailVerified}, UI \\
E2E-009 & Usuário sem papel aceita e a sessão publica \texttt{Aluno} sem
  relogin & \texttt{reconciliarClaimsUsuario}, \texttt{getIdTokenResult(true)},
  AuthContext \\
E2E-010 & Recusa via notificação libera a pendência e marca como lida &
  callable \texttt{rejeitarConviteAluno}, \texttt{Chaves\_Unicas} \\
E2E-CAP-1/2 & Turma cheia sem exceção falha; com exceção justificada registra &
  validação de capacidade, payload real do navegador \\
\bottomrule
\end{longtable}

\subsection{Regressões recentes cobertas}

\begin{itemize}
  \item \texttt{228aae5} (inbox restrita ao destinatário) $\rightarrow$ E2E-003.
  \item \texttt{e23ede1}, \texttt{ef4f577}, \texttt{feca32a} (refresh forçado de
    sessão/claims) $\rightarrow$ E2E-009.
  \item \texttt{df6b3da} (teste de refresh forçado) $\rightarrow$ E2E-009.
  \item \texttt{96c25bc}, \texttt{977ee75} (projeção segura de matrícula
    necessária) $\rightarrow$ E2E-007/E2E-008.
  \item \texttt{80a991c} e \texttt{98b4f21} são infraestrutura (escopo de build
    e \texttt{.gitignore}); não exigem E2E.
\end{itemize}

Os testes também observam o \emph{payload real} enviado pelo SDK
(\texttt{postData}) para provar que campos opcionais ausentes não são
serializados como \texttt{null}: \texttt{justificativaExcecao},
\texttt{viaNotificacao} e \texttt{matriculaInformada}. Esse foi o defeito que
originou a rodada.

\subsection{Observabilidade e isolamento}

A guarda de console falha em \texttt{console.error} e \texttt{pageerror} não
esperados. Testes negativos declaram padrões permitidos; a única exceção global
é o aviso de reconexão do SDK Firestore em cold start, que é transporte e não
viola regra. Testes que preferem a via externa recuperam o OOB do Auth Emulator
como transporte do e-mail, nunca fabricam \texttt{tokenConvite}.

\subsection{Divergência de produto registrada}

\textbf{Sintoma.} Uma conta autenticada com e-mail verificado e \emph{nenhum}
papel persistido (fixture \texttt{aluno.authOnly}) recebe a notificação interna
\texttt{CONVITE\_PARA\_TURMA} do backend, mas é redirecionada para
\texttt{/desativado} por \texttt{ProtectedRoute}, de modo que o sino (UI-12) não
é montado e a notificação fica inalcançável pela interface.

\textbf{Fonte normativa.} Seção~8 UI-12 (caixa única por UID, consultada pelo
destinatário autenticado) e Seção~9 ALU-01a (via interna pelo sino), com a
leitura de bootstrap já autorizada pela Rule de M13 (IMP-RULES-003).

\textbf{Causa localizada.} \texttt{frontend/src/components/auth/ProtectedRoute.tsx}
trata \texttt{roles.length === 0} como bloqueio total, impedindo o bootstrap de
convite na UI embora o backend e as Rules já o suportem.

\textbf{Impacto na suíte.} E2E-009 exercita o aceite real
(\texttt{viaNotificacao}) lendo o deep link do próprio documento de notificação
persistido, uma vez que a superfície de UI correspondente não existe para contas
sem papel. O aceite e o refresh de claims são provados; a descoberta da
notificação pela UI permanece não conforme. A correção do gate de navegação
exige decisão de UX (shell autenticado sem papéis) e não foi aplicada nesta
rodada.
